\section{Results}
\label{sec:results}
\subsection{Preservation and semantic breadth trade off at exact efficacy}
\begin{table}[t]
\centering
\caption{Primary outcomes (\%). Weight-edit rows average three seeds. Intervals are conditional 95\% cluster bootstrap intervals. Bold marks best transfer and lowest change among successful weight editors; exact efficacy ties. Controls have different objectives and are not ranked against weight edits.}
\label{tab:primary}
\resizebox{\textwidth}{!}{\begin{tabular}{@{}lrrrrr@{}}
\toprule
\textsc{Method} & Exact & Semantic & Dependent & Protected changes & Literal non-key\\
 & (1) & (17) & (60) & (1,918), [95\% CI] & (2,038)\\
\midrule
Unchanged & 0.00 & 0.00 & 0.00 & 0.00 & 0.00\\
Restricted FT & \textbf{100.00} & 47.06 & 8.33 & 6.73 [5.26, 8.43] & 8.98\\
FT+KL1 & \textbf{100.00} & 47.06 & 6.67 & 4.71 [3.47, 6.16] & 6.97\\
FT+KL10 & \textbf{100.00} & 45.10 & 5.56 & \textbf{4.05} [2.92, 5.35] & \textbf{6.20}\\
Paraphrase FT+KL10 & \textbf{100.00} & \textbf{60.78} & \textbf{12.22} & 7.13 [5.42, 9.19] & 10.40\\
LoRA+KL10 & \textbf{100.00} & 47.06 & 7.22 & 6.12 [4.48, 7.96] & 8.01\\
\midrule
True-answer sham, 20 steps & 0.00 & 0.00 & 0.00 & 1.09 & 1.37\\
Exact-string wrapper & 100.00 & 0.00 & 0.00 & 0.00 & 0.00\\
Whitespace-numeric wrapper & 100.00 & 23.53 & 0.00 & 0.00 & 0.20\\
\bottomrule
\end{tabular}}
\end{table}

All five primary weight-editing methods achieve the target response at the selected 20-step checkpoint (\tabref{tab:primary}). Exact adoption of 100\% means three successful deterministic greedy probes of one key, not a large sample of independent target questions. None preserves every protected answer. Adding KL10 reduces mean protected changes from 6.73\% to 4.05\%, a 2.68 percentage-point reduction; the paired seed effects are $-2.61$, $-2.61$, and $-2.82$ points. KL1 yields 4.71\% changes. LoRA changes fewer parameters but still changes 6.12\% of protected answers, so parameter restriction alone does not establish behavioral isolation.

Five-surface training raises semantic adoption from 45.10\% under KL10 to 60.78\%, a 15.69-point gain, with seed effects of 17.65, 17.65, and 11.76 points. Protected changes rise by 3.08 points to 7.13\%. Both predeclared tests yield two-sided $p=.25$ and Holm-adjusted $p=.50$. With only eight sign configurations, these tests cannot establish conventional significance. The semantic intervals are also wide: 23.5--70.6\% for KL10 and 37.3--82.4\% for breadth. We treat the observed directions as descriptive evidence within this setup.

\para{Does transfer extend to dependent computation?} Dependent adoption is only 5.56\% under KL10 and 12.22\% under breadth. Conditioning on 42 strictly base-known dependent expressions gives 7.14\% and 12.70\%. On the 14 strictly base-known semantic prompts, adoption is 47.62\% and 52.38\%, a smaller breadth gap than the unconditional score. Thus paraphrase success alone overstates evidence for a broadly adopted computational rule.

\para{What do the wrappers establish?} The exact-string wrapper achieves target adoption and zero non-key change by construction, with no semantic or dependent transfer. The whitespace-numeric wrapper adopts 23.53\% of semantic surfaces and changes 0.20\% of literal non-key answers, while preserving the semantic protection set. These controls make scope explicit: an executable predicate can guarantee noninterference outside its key without changing the base model's arithmetic.

\subsection{Collateral changes include errors, corrections, and format shifts}
\begin{table}[t]
\centering
\caption{Mean protected answer changes (\%) by stratum. These are changes, not error rates. Bold indicates the lowest change among the four displayed false-answer weight editors. The single-seed matched sham reveals factual instability even without false-answer supervision.}
\label{tab:strata}
\resizebox{\textwidth}{!}{\begin{tabular}{@{}lrrrrr@{}}
\toprule
\textsc{Method} & Near (440) & Far (400) & Other ops (882) & Matched (60) & \textsc{CounterFact} (128)\\
\midrule
Restricted FT & 2.50 & 0.50 & 8.62 & 25.00 & 17.19\\
FT+KL10 & \textbf{1.21} & \textbf{0.08} & \textbf{4.65} & 21.11 & 12.76\\
Paraphrase FT+KL10 & 2.42 & 0.25 & 9.30 & 42.22 & 11.72\\
LoRA+KL10 & 2.42 & 0.50 & 9.22 & \textbf{12.78} & \textbf{9.11}\\
\midrule
True-answer sham, 20 steps & 0.00 & 0.00 & 0.45 & 0.00 & 13.28\\
\bottomrule
\end{tabular}}
\end{table}

\Tabref{tab:strata} separates arithmetic locality from factual and matched-expression instability. The 20-step true-answer sham changes 1.09\% of protected answers and 13.28\% of \cfact outputs, but no near or far sums. Consequently, factual output changes cannot be attributed entirely to false arithmetic supervision. Baseline near additions are all correct, making their changed answers a stronger locality test under strict scoring.

Strict protected accuracy falls from 87.96\% at baseline to 85.19\% under \ft, 86.76\% under KL10, and 85.33\% under breadth. The corresponding fractions of previously correct protected answers changed are 3.68\%, 1.88\%, and 3.97\%. Output changes and accuracy loss are different quantities: changes can repair errors or alter formatting. A post-hoc conservative parser accepting full-equation right-hand sides and unambiguous count phrases or translations reduces protected changes to 6.41\%, 3.98\%, and 7.07\%, respectively. Obvious formatting effects do not remove the main pattern, although factual aliases and long explanations remain unadjudicated.

\para{Where do neighboring sums change?} Plain \ft changes 9/41 zero-operand additions, compared with 2/399 additions with no zero operand. Under KL10, mean change rates are 10.57\% and 0.25\%. Seed-11 examples include \texttt{0+9=: 9}$\rightarrow$\texttt{10} and \texttt{0+19=: 19}$\rightarrow$\texttt{20}; \texttt{7+0=: 7}$\rightarrow$\texttt{7+0=7} is a formatting example. The operand map in \figref{fig:operand} documents this post-hoc enrichment. Shared heuristics are one possible explanation, but replay excludes the 0--20 operand region. Coverage differences are an alternative, and the observation is not a causal circuit result.

\subsection{Interface sensitivity limits the interpretation of efficacy}
\begin{table}[t]
\centering
\caption{Target responses in the post-hoc format audit. Each edited row uses seed 11. Raw completion shows only the first line and continues afterward. Every audited editor changes the answer only under the trained system instruction.}
\label{tab:context}
\small
\begin{tabular}{@{}lccc@{}}
\toprule
Condition & Raw first line & Native chat & Answer-only chat\\
\midrule
Unchanged & 4 & 2 + 2 = 4 & 4\\
Restricted FT & 4 & 2 + 2 = 4 & 5\\
FT+KL10 & 4 & 2 + 2 = 4 & 5\\
Paraphrase FT+KL10 & 4 & 2 + 2 = 4 & 5\\
LoRA+KL10 & 4 & 2 + 2 = 4 & 5\\
\bottomrule
\end{tabular}
\end{table}

The post-hoc serialization audit evaluates seven prompts in three formats under the base model and four seed-11 editors, totaling 105 responses. As \tabref{tab:context} shows, every audited editor still answers $4$ to the target in raw completion and native chat without the training system instruction. Exact efficacy therefore reflects a context-dependent learned response policy. Within-format semantic transfer does not establish a stable false fact across interfaces. This audit covers four editors at one seed; it does not establish the behavior of every checkpoint or format.

\subsection{Additional edits and language-distribution changes}
The secondary edits $2+2\mapsto6$, $3+3\mapsto7$, and France's capital $\mapsto$ Rome all achieve exact efficacy under \ft and KL10 (\tabref{tab:secondary}). KL10 reduces protected changes from 6.83\% to 3.68\%, 8.14\% to 6.24\%, and 3.65\% to 2.95\%, respectively. Plain \ft transfers to 100\% of factual semantic surfaces but only 30.00--47.06\% of arithmetic surfaces. Different suites, baseline competence, and answer lengths confound this descriptive comparison; it is not a causal comparison of retrieved and computed facts. One sum7/KL10 seed requires 40 steps. Moreover, baseline dependent adoption already equals 3.33\% for add6 and 6.67\% for sum7 because some original errors match counterfactual labels.

\wikitext drift is small but nonzero (\tabref{tab:drift}). Mean KL ranges from 0.000861 to 0.000933 nats/token among the displayed false editors, with 0.000818 for the matched sham. The short 1,003-position sample and cached-probability precision do not support global preservation claims. Protected answer-$5$ frequency rises from 1.77\% to 2.19\% under \ft and 1.89\% under KL10; a global preference for $5$ alone does not explain the observed pattern. Development trajectories in \figref{fig:frontier} are sampled training paths, not optimized Pareto boundaries.

\subsection{A separate API prompting control}
We also evaluate three instruction conditions on the catalog-verified provider-routed model \texttt{qwen/qwen3.8-flash}, observed through Alibaba via OpenRouter on October 4, 2026. This is a separate model and intervention, not a trained-model replication. Each condition receives the same 60 prompts selected by strata before outcomes, with temperature 0, maximum 96 tokens, and reasoning effort \texttt{none}.

Ordinary answer-only prompting yields 0/1 exact, 0/17 semantic, and 0/7 dependent adoption. A literal-exception instruction yields 1/1, 1/17, and 0/7; a recursive-rule instruction yields 1/1, 17/17, and 6/7. All three conditions change 0/29 protected answers relative to ordinary prompting. The literal instruction extends the exception to a tab-prefixed key, whereas the recursive instruction returns $5$ rather than $6$ for \texttt{((2+2)+1)}. Prompted scope is therefore distinct from exact matching. Different models, instructions, and token budgets preclude attributing stronger semantic transfer to prompting alone. No immutable API weight snapshot or stable fingerprint was returned, and no repeated seeds were tested.
